% Six-part dossier template
% Compile main.tex with pdfLaTeX (run twice for the table of contents).
\documentclass[11pt,a4paper,oneside]{report}

% ---------- Packages ----------
\usepackage[utf8]{inputenc}
\usepackage[T1]{fontenc}
\usepackage[english]{babel} % Change to [swedish] if appropriate
\usepackage{lmodern}
\usepackage[a4paper,margin=27mm,headheight=15pt]{geometry}
\usepackage{microtype}
\usepackage{graphicx}
\usepackage{xcolor}
\usepackage{titlesec}
\usepackage{fancyhdr}
\usepackage{enumitem}
\usepackage{booktabs}
\usepackage[hidelinks]{hyperref}


% ---------- Math Packages ----------
\usepackage{amsmath, amsfonts}

% ---------- Dossier information: edit these ----------
\newcommand{\DossierTitle}{Modern Experimental Mechanics}
\newcommand{\DossierSubtitle}{Lecture Dossier}
\newcommand{\AuthorName}{Vilmer Dahlberg}
\newcommand{\Institution}{Division of Solid Mechanics, Lund University}
\newcommand{\DossierDate}{\today}

% ---------- Appearance ----------
\definecolor{DossierInk}{HTML}{203047}
\definecolor{DossierAccent}{HTML}{287C83}
\definecolor{DossierMuted}{HTML}{64748B}
\hypersetup{colorlinks=true,linkcolor=DossierAccent,citecolor=DossierAccent,urlcolor=DossierAccent,
  pdftitle={\DossierTitle},pdfauthor={\AuthorName}}
\titleformat{\chapter}[display]
  {\normalfont\bfseries\color{DossierInk}}
  {\filleft\Large\chaptername\ \thechapter}
  {0.8ex}
  {\titlerule\vspace{1.2ex}\Huge}
  [\vspace{0.8ex}\titlerule]
\titlespacing*{\chapter}{0pt}{-12pt}{28pt}
\setlist{itemsep=0.25em,topsep=0.4em}
\setlength{\parindent}{0pt}
\setlength{\parskip}{0.65em}
\pagestyle{fancy}
\fancyhf{}
\fancyhead[L]{\small\color{DossierMuted}\nouppercase{\leftmark}}
\fancyhead[R]{\small\color{DossierMuted}\DossierTitle}
\fancyfoot[C]{\small\color{DossierMuted}\thepage}
\renewcommand{\headrulewidth}{0.4pt}
\renewcommand{\headrule}{\hbox to\headwidth{\color{DossierAccent}\leaders\hrule height \headrulewidth\hfill}}

\begin{document}

% ---------- Cover ----------
\begin{titlepage}
  \thispagestyle{empty}
  \vspace*{22mm}
  {\color{DossierAccent}\rule{\textwidth}{2pt}}\\[14mm]
  {\Huge\bfseries\color{DossierInk}\DossierTitle\par}
  \vspace{5mm}
  {\Large\color{DossierMuted}\DossierSubtitle\par}
  \vfill
  {\large\bfseries\AuthorName\par}
  \vspace{2mm}
  {\Institution\par}
  \vspace{2mm}
  {\DossierDate\par}
  \vspace{14mm}
  {\color{DossierAccent}\rule{\textwidth}{2pt}}
\end{titlepage}

\pagenumbering{roman}
\chapter*{Preface}
% Replace this note with a short overview of the dossier, its purpose, and how
% the six parts fit together. Delete this page if you do not need a preface.
This dossier is organized into six parts - one for each subject in the course. 

\section*{Generative AI Disclaimer}
The use of generative AI was restricted to spell checking and \LaTeX file formatting.

\addcontentsline{toc}{chapter}{Preface}

\tableofcontents
\clearpage
\pagenumbering{arabic}

% ---------- Six parts ----------
% Keep each part in its own file. Change titles in the matching part file.

% Experimental method in mechanics and full-field methods
\chapter{Experimental method in mechanics and full-field methods}

\subsection{Experimental mechanics and the experimental method}
%• What is an experiment?
%• Why do we do experiments?
In my opinion, the definition of an experiment is very simple: An experiment is a procedure in which we seek to learn more about a process with one or more inputs and outputs. And in experimental mechanics, the processes we seek to understand are mechanical in nature. There are many reasons to do experiments. Because it's fun. Because we want to create models of reality to be used for engineering, or otherwise.
 
Now this definition not constitute a \textit{good} experiment per say. 

Broadly speaking, in experimental mechanics we seek a deeper understanding of the mechanical behavior of mechanical systems and materials. I don't exactly know what that means but something like that. We are also doing some experiments, yes, experiments- This is because we want to understand the mechanical behavior. Well I guess that is just saying the same thing twice but hey I need to fill this out so I can reach at least two pages. So we are doing experiments, and doing them in a controlled manner, so as we can eliminate spurious external influences and just look at one thing at the same time (there is a word for this i just cant think of it) and so we define dependent variables and independent variables. and we vary the independent variables (or perhaps just one at a time) and then observe the dependent variables. Hence we need to be able to make observations of our systems.


%• What do we want to measure in experimental mechanics?
Historically we have only been able to observe macroscopic behaviors, such as the force or displacement. But in modern times we are able to to full-field analyses.

%• How might we use the information we gain?
So after we've made some observations of our system we construct models which describe this behavior, so we are able to make predictions of other systems which are similar but not identical. This is in some sense the back-bone of engineering, which requires precise models to be able to construct great things.

Often times, even though we control our experiment meticulously there will be some spurious or random processes, which is why we often want to repeat the experiment multiple times. We therefore \textit{want} to generate more data than our model requires. Thus when analyzing the data we have to do some kind of least squares or better yet take randomness and uncertainties (so maybe here i am looking for the word standard deviation) into account when fitting the data to our model. If we don't collect enough data we risk over-fitting the data to our model, in which case we might have very bad predictions outside the range of data we have collected.

You also have to be careful and think twice about what we're actually measuring. For example, to measure a force we often use a load cell, which in reality actually measures the stretch of a strain gauge, which actually measures the resistance of a wire being stretched which is actually measuring a voltage, and I don't know if i can go deeper here but you get my idea. We have to be careful to make sure that what we're measuring is within the range of the instrument we're using (i guess) and other things. Sometimes it can be hard to know exactly what we're measuring (example with the experiment where the change in temperature changed the strain and so on).

You also have to be careful about what are facts and what are inferences you can make based on facts. You can also make judgments, which are subjective observations. They reveal our values and feelings towards something (this is copied verbatim).

%• What are resolution, accuracy and precision?
Given a constant process, i.e. experiment, and a measuring device \textit{accuracy} describes how consistent multiple measurements are no matter how close it is to the expected value. I.e. a high accuracy measuring device has low standard deviation. On the other hand, precision describes how close multiple measurements are to the true value. I.e. a high precision measuring device has a low deviation from the expected value of the process. All observations which are made of the world are at the end of the day digital, i.e we must produce a single value for each measurement. The resolution of a measuring device is the smallest amount two different results can deviate.


\subsection{Full-Field and in-situ measurements in experimental mechanics}
%• What are full-field measurements and why are they of interest in experimental mechanics?
As mentioned above, historically we've only been able to do measurements of macroscopic (global) quantities. This is because measuring local quantities have historically required destructive methods. Which makes sense, how else would you look "inside" something? Unfortunately, it's very difficult to draw any conclusions from macroscopic experiments which can be applied to other systems. Often times we seek some kind of "local" material model, so that we can apply the knowledge to other "geometries". This requires "full field" measurements. 

Many modern materials are heterogeneous, and understanding their mechanical properties requires not only full-field analyses, but also multiple scales. Sometimes we also want to know multi-modal knowledge, for example paper-based materials are very sensitive to moisture and temperature. Small changes can have big effects on the mechanical response.

%• What sorts of fields can be measured?
There are many types of fields we would like to measure, for example the orientation of fibres in wood/paper materials. The grain orientation in a metal or granular structure. The chemical expansion of a battery during charging/decharging. Tracking dislocations (plasticity) in loading. Crack growth. Phase transformations. Material damage and fatigue.

%• What is meant by “in-situ” in the context of experimental mechanics?
In-situ experiments are experiments done during some kind of mechanical loading or otherwise. I.e. the independent variable may be a load or flow of something. We are in other words not looking at something static, or something \textit{postmortem}, i.e. after the deformation.

%• Why are in-situ experiments of interest?
We want to do in-situ experiments so that we can separate the geometry from the material/local behavior. In this way we can predict the response of structures consisting of the same material, but with a different geometry. This is a corner-stone of FEM.


% Digital image correlation for Experimental Mechanics
\chapter{Digital Image Correlation}

% Replace this example text with the content of Part 2.
\section{Overview}
Write a brief introduction to this part and explain how it relates to the overall dossier.

\section{Main discussion}
Organize the material with sections and subsections as needed. You can include
figures, tables, equations, and citations here.

%Provide a short overview of DIC:
%• What are the different types of DIC?

%• What are the input data and how are they acquired?

%• What can one measure using DIC (and how does this differ with 2D, 3D-surface and 3D-volume/DVC)?

%• What is local vs global DIC and cumulative versus incremental?
In cumulative DIC we always compare the current image to the first image. In contrast in incremental DIC we compare to the image one or a few time steps back.

In local DIC the transformation of each window is treated independently. In contrast, in global DIC we consider the displacement of each node as a part of a larger mesh of points and solve for the displacements of each point simultaneously.

Local DIC is faster and has less assumptions. For example fracture can be tracked. On the other hand we can get poor correlation in between regions.

Global DIC is a bit slower, and requires global convergence. It also is a bit more constrained, as it cannot capture fracture. On the other hand you get less oscillations. Also can be used to do inverse analysis, meaning you can compute the initial image given any image. 

%• What are the experimental constraints for the different DIC (2D-surface. 3D-surface, 3D-Volume) approaches, i.e., what is required use 2D/3D/Volume DIC to monitor an experiment? (sample preparation, type of deformation, type of images, experimental set-up constraints, etc.)
With one camera, you can only realistically measure what is happening in the plane. For 3D-surface DIC we require calibration. 

Digital volume correlation compares digital volumes, requires radiography. The sequence is deform, acquire image, etc.

In projection-based DVS we can somehow use predictions of the deformation. 

If you have two cameras you can detect depth, i.e. now we can measure either out-of-plane displacement or measure non-plane surfaces.

With tomography we're able to compute the full strain tensor.

%• Describe a basic (local) DIC analysis processing flow (DIC algorithm)
The basic DIC analysis processing flow is as follows. 

%\begin{list}
%	\item[i] Load two images to compare\\ %$\mathcal{A}$, $\mathcal{B}$ \\
%	\item[ii] Define analysis window\\ %$\mathcal{W} \subset \mathcal{A}$  \\
%	\item[iii] Define motif size\\ %$L$ \\
%	\item[iv] For each pixel\\ %$p \in \mathcal{W}$ \\
%\end{list}

%• Why is subpixel resolution on displacements needed and how can it be achieved
Since the displacement field is rarely integer values the calculated displacement should reflect that, i.e. be floats. Otherwise we risk getting checkerboarding.

When we then compute the strain, we require gradients of the displacement field. If we then only have integer displacements we can get very large jumps. 

%• Why might one use DIC in experimental mechanics?
Identify boundary conditions, rather than assuming that we know what they are. 

Identify non-homogenous phenomena

Identify model parameters, to enrich finite element analysis. 

Identify boundary conditions.

Virtual fields methods.

Integrated digital image correlation..

%• How can DIC be used to develop improved material models


% Acoustics, Ultrasonics & Ultrasonic tomography for Experimental Mechanics
\chapter{Acoustic and Ultrasonic Methods}

% Provide a summary of elastic wave properties/characteristics

% What are the different types of waves and how do they differ in their propagation?

A wave is not a phenomena which exists in its own right, rather it emerges from time or space dependent systems or fields. Thus waves can only exist in some media, be it a solid (as is the case in experimental mechancis) or fluid, or even electro-magnetic waves such as light. Since waves are an emergent phenomena, a wave can be anything. The movement of tightly packed humans (such as at concerts) is well-approximated by a plane fluid. In the context of experimental mechanics we are interested in waves in solid mater.

Waves are often times sinusoidal, but can in principle be any shape. For instance surface waves (Rayleigh and Love waves), or the the wavefronts of a circular membrane (such as a drum) are not sinusoidal in the radial direction. Worth noting is that in space-time systems, standing waves, which are fixed in space but time variant, can emerge. Non-sinusoidal waves can appear even for regular geometries (i.e. infinite geometries). However, most systems have some amount of energy loss. This loss is often times proportional to the wave frequency (higher frequencies having smaller lifetimes than lower). As such, non-sinusoidal waves, such as square waves (which contain lots of high-frequency components) often degenerate to their lowest frequency components, which (depending on the geometry) are sinusoidal.

% What is an elastic wave?
In the context of experimental mechanics we are interested in waves propagating in solids. The governing equation is given by the balance of linear momentum, which is written in strong form as
\begin{align}
	\nabla \cdot \boldsymbol{\sigma} + \boldsymbol{b} 	& = \rho \ddot{\boldsymbol{u}} 	& \text{in } & \Omega \label{eq:linear_momentum}\\
	\boldsymbol{\sigma}\cdot\boldsymbol{n}				& = \overline{\boldsymbol{t}}				& \text{on } & \partial\Omega^t
\end{align}
where $\boldsymbol{\sigma}$ is the Cauchy stress, $\nabla$ denotes the gradient, $\boldsymbol{b}$ are the body forces, $\rho$ is the material density and $\boldsymbol{u}$ are the displacements. $\Omega$ denotes the body of interest and $\partial\Omega^t$ the portion of its surface subject to traction $\overline{\boldsymbol{t}}$. For free waves, we assume zero body forces $\boldsymbol{b} = \boldsymbol{0}$. Furthermore, for purely elastic waves, we assume a linear constiutive relationship, $\boldsymbol{\sigma}=\mathbb{D}\mathrel{:}\boldsymbol{\epsilon}$. Note that the consitituive matrix $\mathbb{D}$ may vary spatially. Additionally, while this restricts our discussion to small strains, it's worth noting that even large-strain elastic waves are possible.

The solution of \eqref{eq:linear_momentum} depends on which body is studied and the boundary conditions. Unfortunately, its not possible to find a general solution. There are however some special cases. For instance, under plane strain conditions the solution is given by Rayleigh and Love waves \textcolor{red}{refs}. For the elastic waves in a string, the wave equation becomes
\begin{equation}
	\ddot{u} = v^2 \partial^2_x u
\end{equation}
where $u$ now only has one component, and is the string's normal deflection. The characteristic solutions to the above are given by $u(t, x) = u_0(x \pm vt)$ where $u_0$ is the initial condition and $v$ is the wave-speed. In 3D media, a similar wave equation can be found, except there are now $3$ individual components, and for each component of $u$ there for all three dimensions. Many of these are the same if the material is well approximated by isotropy, but in general the wave speed in any component is different depending on the wave's direction of propagation. \textcolor{red}{cite and find expression for this.}.

% How can we transmit/receive elastic waves? What are acoustic emissions and how are they measured? How might the information be used to gain understanding about the mechanisms of deformation in different materials?

Elastic waves in solid matter can be detected with Piezoelectric contact transducers which use precisely tuned voltage-sensitive quartz crystals (maybe ref here). Importantly, these transducers can be used to both measure and produce waves. As such, experiments can be passive, wherein a transducer is attached to the specimen and waves are induced by the deformation of the specimen. Experiments can also be active, where the specimen is actuated using the transducer and the reflection of the signal is measured with the same transducer.

Examples of passive acoustic experiments can be found in....

For example, the acoustic emissions of a specimen can be used to infer the creation or propagation of crack. Another example is geology, a signal is induced by the earth's crust or third party human sources (read: bombs).

Multiple transducers can be placed on a single specimen. The measured acoustic signal can be then used to estimate the origin of an event, given some model of the wave propagation. It is also possible to characterize the polarity of an event (i.e. shear, tensile  or compressive) using multiple sensors placed on either end of a symmetric specimen. This is known as the polarity method \textcolor{red}{cite}. 

% What can be measured in terms of data from recorded elastic-wave waveforms?
For a single source we can perform a fourier transformation, which exposes the amplitude and frequency information of the signal. If the source of an event is known, we can correlate it with the measured signal to give an estimate of the material' wave speed. Or, if multiple transducers are mounted to a single sample we can triangulate the location of an event by assuming that the wave eminating from the event source travels elastically and at constant speed. 

% What information can be gained about a material from measurements of elastic wave velocities?
Can be used to describe the mechanisms of failure.

Something about wood cracks?

Wood example: there is constant brittle deformation (whatever that means).

Tetrapak, used to characterize brittle failure.

Porosity

Damage

% What is anisotropy (the definition and also the material causes of it) and how might it be measured with ultrasonic measurements?
Virtually all materials behave differently depending on their orientation in space. For instance paperboard, which consists of paper fibers oriented in a principle direction, is stiffer when stretched in the direction of the fiber orientation, and weaker when pulled in the orthogonal direction. This is known as anisotropy, meaning \text{not the same in all directions}. The degree of anisotropy can, however, vary widely, paperboard being on the far end of high anisotropy. Some materials are however almost identical in all orientations, and are treated as isotropic, meaning \textit{the same in all directions}. This information is key for developing efficient material models, which are greatly simplified if the material is well approximated by isotropy. For example, metals are very well approximated by isotropy. Whether or not a material is anisotropic or isotropic depends on its microstructure.  

To test if a material is isotropic we can try to measure the wave speed in different orientations. If they are the same in all principle directions, the material is at least orthotropic. If it's completely independent on the spatial orientation the material is isotropic. On the other hand, we can, for instance, estimate the (edward and thomas) fiber direction of a material by measuring the wave speed in a number of directions. The fiber direction should then have the fastest wave-speed, and vice versa. 

Extending on this idea, some materials behave hyper-elastically, for example rubber which has a wide elastic range. The material properties, for instance the wave speed, are in any point of these materials dependent on the stress state at that point. 

% Reflection and refraction
In the transition between two material phases with different elastic properties.

% What is ultrasonic tomography? What data are needed to perform ultrasonic tomography and what is the procedure for analysing the data to produce a map of ultrasonic velocity (or attenuation)

The goal of ultrasonic tomography is to make full-field measurements of a specimen by measuring the acoustic signal at many points along its boundary. By introducing a model of the specimen's behavior, model parameters can be fit using the gathered data. For example, the specimen can be discretized into any number of elements, each of which has it's own wave speed. The spatial resolution is bounded from below by the signal's largest wavelength in the specimen. Since the wave may refract and bend where the wave speed changes, the wave speed at any point will be dependent on the signal intensity at that point, and hence the model is non-linear. Fortunately, this problem is only weakly non-linear, so an iterative procedure is expected to converge quickly. (Hall et al 2005, 2007).



% X-ray scattering for Experimental Mechanics
\chapter{Part 4 Title}

% Replace this example text with the content of Part 4.
\section{Overview}
Write a brief introduction to this part and explain how it relates to the overall dossier.

\section{Main discussion}
Organize the material with sections and subsections as needed. You can include
figures, tables, equations, and citations here.

% What are x-rays (e.g., in terms of wavelength and type of radiation)
X-rays are highly-energetic wavelets of electromagnetic radiation (light). The wavelength is between ultraviolet light and gamma-rays, roughly in the range of $10^{-10\pm 2}$. The attenuation of x-rays is dependent on the material. This can be used to work out the type of material. 

% How are x-rays produced and what are the key parameters defining the nature of an x-ray source and how do different sources differ in their mechanisms and parameters
1. Lab source
	easy to use but poor resolution
	
2. Synchrotron
	better resolution but restricted availability
	

	
X-rays can be produced by exciting electrons, and allowing them to relax thereby producing a photon with a predictable energy (given by the difference in energy level between the states).

X-rays can also be produced by accelerating electrons (using for instance a linear accelerator), and then changing their trajectory by forcing them around a circle, which produces light. The electrons are accelerated into the storage storage ring, which are steered into a circle by the use of strong magnets.

They can also be accelerated into a target, which produces light via the nuclear interaction between the electron and target. The key parameters which determine the energy of the photons is the voltage of the 

% How do x-rays interact with matter and which are the important mechanisms for experimental mechanics? What makes x-rays a good probe for experimental mechanics?
The higher the energy, the larger the cross-section.

The convergance of the beam is very important for the detail of the imaging. If the beam is very convergent, we can infer from the divergance of the image some information about the specimen.

High flux gives fast imaging.

% If you wanted to study hard matter (e.g., metals) what might you need in terms of the x-ray energy? (and conversely what would you need for soft matter, e.g., polymers?)

Wavelength is comparable to distance between crystals, so we can use diffraction.

High-photon energy ionizes atoms, so they modify the material (destructive). 

X-rays interact both magnetically and electrically, but the magnetic interaction is very weak.

There are five typical interactions, of which the first three in the list are the most intersting for x-ray imaging in mechanics.

Incoherent scattering not too interesting for us, as it's random? What does that mean.

Coherent scattering is useful since its elastic, and we can predict the result.



% X-ray tomography for Experimental Mechanics
\chapter{Part 5 Title}



% Neutrons for Experimental Mechanics
\chapter{Part 6 Title}

% Replace this example text with the content of Part 6.
\section{Overview}
Write a brief introduction to this part and explain how it relates to the overall dossier.

\section{Main discussion}
Organize the material with sections and subsections as needed. You can include
figures, tables, equations, and citations here.

% Example figure:
% \begin{figure}[htbp]
%   \centering
%   \includegraphics[width=0.75\textwidth]{figures/example.pdf}
%   \caption{A descriptive caption.}
%   \label{fig:example}
% \end{figure}

% Example table:
% \begin{table}[htbp]
%   \centering
%   \caption{A concise table caption.}
%   \begin{tabular}{ll}
%     \toprule
%     Item & Description \\
%     \midrule
%     Example & Replace with your content \\
%     \bottomrule
%   \end{tabular}
% \end{table}


% Optional bibliography: uncomment if you use BibTeX.
% Add citations in the text with \cite{key}, then compile with BibTeX.
% \bibliographystyle{plain}
% \bibliography{references}

\end{document}
